% Paper A -- Pattern Recognition submission format (Guide for Authors, checked 2026-11-22):
% single column, 10 pt Times, double spacing, margins 4.3/4.8/4.3/4.8 cm (T/R/B/L), 20-35 pages incl. refs.
% Proofs and further results are in the separate supplement.tex (xr cross-references; compile main, supplement, main).
\documentclass[preprint,10pt]{elsarticle}
\usepackage[top=4.3cm,right=4.8cm,bottom=4.3cm,left=4.8cm]{geometry}
\usepackage{times}
\usepackage{amsmath,amssymb,amsthm,booktabs,graphicx,array,longtable}
% figures are drawn at 17 cm (PR full-page width) and may extend into the wide review margins
\newcommand{\widefig}[1]{\makebox[\linewidth][c]{\includegraphics[width=17cm]{#1}}}
\usepackage{etoolbox}
\usepackage{xr-hyper}
\usepackage[colorlinks=true,allcolors=blue]{hyperref}
\externaldocument[][nocite]{build/supplement}[Paper_A_supplement_EN.pdf]
\linespread{1.667}% equals setspace's \doublespacing at 10 pt
% let floats share pages with text instead of leaving half-empty float pages
\renewcommand{\topfraction}{0.9}\renewcommand{\bottomfraction}{0.8}\renewcommand{\textfraction}{0.07}\renewcommand{\floatpagefraction}{0.85}\setcounter{topnumber}{2}\setcounter{totalnumber}{3}
\newtheorem{proposition}{Proposition}
\newtheorem{corollary}{Corollary}
\journal{Pattern Recognition}
% title 14 pt and author names 8 pt, as the PR guide and editorial office require for the title page
\makeatletter\patchcmd{\pprintMaketitle}{\Large\@title}{\fontsize{14}{17}\selectfont\@title}{}{\typeout{PATCHFAIL title}}\patchcmd{\pprintMaketitle}{\normalsize\elsauthors}{\footnotesize\elsauthors}{}{\typeout{PATCHFAIL authors}}\makeatother
\pdfsuppressptexinfo=-1% do not copy the producer metadata of included figure files
\begin{document}
\begin{frontmatter}
\title{When the site predicts the diagnosis: evaluation, bridging and within-site training for multi-site neuroimaging classifiers}
\author[ee]{Jincheng Hu}
\author[ee,klicn,hebei]{Ping Xie}
\author[ee]{Sha Wang}
\author[ee]{Jiaxin Li}
\author[ee,klicn]{Guoqian Jiang\corref{cor1}}
\ead{jiangguoqian@ysu.edu.cn}
\cortext[cor1]{Corresponding author.}
\affiliation[ee]{organization={School of Electrical Engineering, Yanshan University}, city={Qinhuangdao}, postcode={066000}, country={China}}
\affiliation[klicn]{organization={Key Laboratory of Intelligent Control and Neural Information Processing, Ministry of Education, Yanshan University}, city={Qinhuangdao}, postcode={066000}, country={China}}
\affiliation[hebei]{organization={Hebei Key Laboratory of Intelligent Rehabilitation and Neuroregulation, Yanshan University}, city={Qinhuangdao}, postcode={066000}, country={China}}
\begin{abstract}
Diagnostic models for brain disorders are often trained on data pooled across sites. In clinical studies, a site often contributes a single diagnosis, for example patients from a clinic and controls from a university laboratory. In such class-pure data the site predicts the label, so a model can use the site as a cue to the diagnosis, and subject-level cross-validation does not detect it. With an additive site-effect model, we show that class-pure data cannot separate diagnosis from site without an untestable assumption, and that adding subjects to the same sites does not reduce the site part of the error. We derive three rules: evaluate by holding out whole sites, recruit a few bridging participants of the missing diagnosis at each site, and train only on within-site patient and control pairs. We tested each rule with pre-specified criteria on resting-state fMRI from three cohorts covering autism and attention-deficit/hyperactivity disorder. Subject-level cross-validation overstated class-pure models by 0.37 to 0.46 in AUC, and two real class-pure collections lost about 0.30 more than mixed sites. Group-robust training and two published multi-site methods either kept the shortcut or lost the diagnostic signal. Bridges of the missing diagnosis outperformed as many additions of each site's own diagnosis, and a design coefficient computed from the site-by-diagnosis table ranked designs by shortcut size before training (Spearman $\rho\le-0.97$). The within-site loss kept the reverse-set AUC above chance, and with every site bridged its test AUC was within 0.01 of the pooled model.
\end{abstract}
\begin{keyword}
multi-site learning \sep shortcut learning \sep batch effects \sep study design \sep evaluation \sep autism
\end{keyword}
\end{frontmatter}

% Section 1 -- Introduction (rewritten 2026-11-23, condensed the same day). Narrative modelled on DeGrave et al. 2021
% (promise -> data collection creates worst-case confounding -> consequence -> what we show) and Song et al. 2020.
\section{Introduction}\label{sec:intro}

% P1 background
Two neurodevelopmental disorders, autism spectrum disorder (ASD) and attention deficit hyperactivity disorder (ADHD), are diagnosed from behaviour, through interviews and observation that are slow and depend on the examiner. Many studies therefore train machine learning models on functional connectivity to find objective markers, and because a single hospital rarely has enough participants, they pool data across sites, as in the ABIDE and ADHD-200 consortia~\cite{dimartino2014,bellec2017}.

% P2 the general risk of pooling
Pooling has a price. Sites differ in scanners, protocols and recruitment, and these differences can be as large as the effects of the disorders~\cite{yamashita2019}. They are also easy to detect: a classifier identified which of $17$ data sets a brain scan came from with an accuracy above $70\%$, where chance was $6\%$~\cite{wachinger2021}. If the site is related to the diagnosis, a model can use the site as a cue and fail at other sites, known as shortcut learning~\cite{geirhos2020,zech2018}.

% P3 the specific problem
This paper studies the most extreme form of this risk. Suppose that a child psychiatry clinic scans children with autism, while typically developing children are scanned at a university laboratory in another city. Each site then contributes participants of a single diagnosis, so knowing the site means knowing the label; we call such sites \emph{class-pure}. They arise naturally, because patients are recruited where they are treated and controls where they are easiest to reach, and they occur in public consortia: one ABIDE-II collection contains only autistic participants, and one ADHD-200 site only controls. Early COVID-19 chest radiograph collections had the same structure, and models trained on them relied on cues to the data source~\cite{degrave2021}.

% P4 why it matters
Subject-level cross-validation cannot reveal this problem: a held-out patient comes from a site that contributed only patients to training, so recognising the site is enough to classify the patient correctly. In our ABIDE-I experiments, subject-level cross-validation reported an area under the ROC curve (AUC) of $0.94$ for a model that reached only $0.58$ at new sites. Accuracies reported from class-pure data may thus owe a substantial part to the site.

% P5 the gap
Harmonization methods such as ComBat protect the diagnosis as a covariate~\cite{johnson2007,fortin2018}, but at a class-pure site diagnosis and site cannot be told apart, so this protection cannot be applied. Debiasing methods in pattern recognition reweight groups or learn from \emph{bias-conflicting} samples~\cite{sagawa2020,zhang2022cnc,wang2025mdn}; here these would be patients at a control site, and there are none. Exact site invariance~\cite{ganin2016,dinsdale2021,kunda2023} would remove the information that separates the classes~\cite{zhao2019}. In genomics, designs in which batch and group coincide are known to defeat both batch correction and prediction across batches~\cite{song2020,parker2012}, but that work concerns laboratory batches. For diagnostic models trained on many class-pure sites, it remains open how much class purity costs, how such a model should be evaluated, and which data to add and how to train with them.

% P6 what we do
Our method acts on the training data, the training objective and the evaluation, so it can be used with standard score-based classifiers. When every site holds one diagnosis, every comparison between a patient and a control is also a comparison between two sites, and only comparisons within a site are free of site effects. We formalize this with an additive model of site effects and derive three rules (Fig.~\ref{fig:overview}): \emph{evaluate} by holding out whole sites and testing whether the model ranks by site; \emph{design} the data by recruiting a few \emph{bridging} subjects of the missing diagnosis at each site, guided by a coefficient computed from the table of subjects per site and diagnosis; and \emph{estimate} the model only from comparisons within sites. To measure these effects against a known truth, we build class-pure training sets from sites that hold both diagnoses and use the unused participants of the other diagnosis to test whether the model has learned the site. We tested each rule on resting-state fMRI from two disorders in three cohorts, with criteria written down before the corresponding experiments.

% P7 significance
From a pattern recognition perspective, class-pure sites are a spurious correlation without any bias-conflicting samples, and our results show that collecting a few such samples by design can break it.

Our contributions are as follows.
\begin{itemize}
\item \textbf{Analysis.} With an additive site-effect model, we show that class-pure data cannot separate the diagnosis from the site without an untestable assumption, and that adding subjects to the same sites does not reduce the site part of the error. A design coefficient $L$, computed from the site-by-class table, measures how much of it a design retains.
\item \textbf{Evaluation.} Subject-level cross-validation overstated class-pure logistic models by $0.37$ to $0.46$ in AUC across three cohorts, and two published multi-site methods failed in the same setting. Holding out whole sites removed the systematic bias, and a \emph{reverse-set} test on held-out subjects of the opposite diagnosis at the training sites revealed shortcuts that other checks missed.
\item \textbf{Design.} $L$ is a data acquisition criterion that can be computed before any data are collected: it ranked designs by the size of the shortcut, with rank correlations of $-0.970$ to $-0.997$ with the reverse-set AUC. A few bridging subjects of the other class per site curbed the shortcut more than the same number of same-class additions (reverse-set AUC higher by $0.10$ to $0.29$).
\item \textbf{Estimation.} We propose a logistic loss on within-site patient and control pairs, a learning objective whose solution is unchanged by additive site shifts (Proposition~\ref{prop:within}). It kept the reverse-set AUC between $0.59$ and $0.67$ in all $12$ bridged cells of ABIDE-I, and with two-way bridging its test AUC was within $0.01$ of the pooled model; as a linear head on a frozen self-supervised encoder it reached a reverse-set AUC of $0.59$, against at most $0.39$ for deep baselines. On ABIDE-II all assessable pre-specified criteria were met, and on ADHD-200 all assessable criteria except the cost of class purity.
\end{itemize}

\begin{figure}[t]\centering\widefig{figures/fig1_designs.pdf}\caption{Training designs with $k=2$ sites per class. Each column is one site (Site 1 to Site 4), and each figure represents five subjects ($m=20$ subjects per site, $b=5$ bridges). (a) Each site contributes one class. (b) One-way bridging adds controls to the case sites. (c) Two-way bridging adds the missing class at every site. (d) Mixed sites with the same total size as the class-pure design in (a). Dashed boxes mark the reverse set, subjects of the opposite class held out at originally class-pure sites; $L$ is the design coefficient of Eq.~(\ref{eq:L}). Bottom: the three steps of the framework (Section~\ref{sec:framework}).}\label{fig:overview}\end{figure}
% Section 2 -- Related work (condensed 2026-11-23). Each subsection: what the line of work does, what it assumes, where this work differs.
\section{Related work}\label{sec:related}

\subsection{Site effects and harmonization}

Harmonization removes the shift and scaling that each site adds to the features. ComBat~\cite{johnson2007,fortin2018} and its deep variants~\cite{an2025deepresbat} protect biological variables such as the diagnosis by including them as covariates, which works only if the protected variable varies within sites. When groups are unevenly distributed across batches, batch correction can shrink true group differences or create spurious ones~\cite{nygaard2016}, ComBat that protects the diagnosis leaks label information when class proportions differ across sites~\cite{nieto2026}, residual site effects can bias accuracy unless the evaluation accounts for sites~\cite{solanes2021}, and when the covariate distributions of two batches barely overlap, causal methods decline to decide whether a batch effect exists~\cite{bridgeford2025}. Traveling-subject designs separate measurement from recruitment differences~\cite{yamashita2019} but are costly and rare. Class-pure sites are the case of no overlap at all. Our design step creates the missing overlap, and our estimator does not need to estimate site effects.

\subsection{Shortcut learning and debiasing}

Deep models often adopt shortcuts that work on standard benchmarks but fail to transfer~\cite{geirhos2020}; in medical imaging the shortcut is often the hospital or acquisition source~\cite{zech2018,degrave2021}. Remedies include group reweighting~\cite{sagawa2020}, contrastive learning across attributes~\cite{zhang2022cnc}, margin-based emphasis of bias-conflicting samples~\cite{wang2025mdn} and class-wise loss balancing~\cite{zuo2026}. All of them remain well defined on class-pure data, but the data lack a contrast between classes within a site, so the diagnostic effect cannot be separated from the site effects without further assumptions (Proposition~\ref{prop:pure}); contrastive and margin-based methods also need bias-conflicting samples, of which there are none. Domain generalization for classes seen in a single domain~\cite{kaai2025fond} transfers invariance learned from classes that appear in several domains; with class-pure sites, invariance learned across sites of the same class leaves the difference between case sites and control sites (Section~\ref{sec:prop3}). Bridging subjects supply the missing bias-conflicting samples, and our within-site loss uses only pairs that share the spurious attribute and differ in label, the hard negatives of contrastive debiasing~\cite{zhang2022cnc}.

\subsection{Multi-site classification, confounded designs and evaluation}

Methods for multi-site classification from functional connectivity were mostly developed on data in which most sites contain both classes, for example site-aware population graphs~\cite{parisot2018,wu2027croip} and domain adaptation~\cite{song2024braindas,fang2025scda}; federated learning~\cite{wang2025gafd}, minimization of site dependence~\cite{kunda2023} and adversarial or mixed-effects removal of site information~\cite{ganin2016,dinsdale2021,zhao2020cfnet,nguyen2023armed} were also developed and evaluated on such data, as was CORE~\cite{wang2026core}, which corrects confounders site by site. We test two of them in Section~\ref{sec:res-pureonly}. In genomics, confounding between batch and group inflates cross-validated performance~\cite{soneson2014}, and when each of two batches held one outcome, a predictor trained on them failed on the unseen combinations of batch and outcome~\cite{parker2012}. Design-based remedies place shared material in several batches: reference panels~\cite{song2020}, remeasured controls~\cite{ye2023} and technical bridging samples~\cite{bridge2022}. Our bridging subjects are instead new individuals of the other class at the same site, with no repeated measurement, so whether they work depends on whether the site effect is additive. For evaluation, leaving out whole sites is an established protocol in multi-site autism classification~\cite{abraham2017}. Our leave-site-pair-out evaluation applies leave-pair-out cross-validation, which gives nearly unbiased AUC estimates in small samples~\cite{airola2011}, to whole sites; because discrimination within clusters differs from discrimination between clusters~\cite{vanklaveren2014}, Section~\ref{sec:eval} states which of the two it estimates. Evaluation protocols can also hide leakage through model selection~\cite{yang2027} or instability across training domains~\cite{zhang2025pf2hc}.

% Section 3 -- Problem setting and theory. Full proofs: Supplementary Section S1 (sections/supp_a.tex).
\section{Problem setting and theory}\label{sec:theory}

With a deliberately simple model, in which each site adds the same shift to the features of all its participants, we ask whether class-pure data determine the diagnostic effect at all (Section~\ref{sec:prop1}), how much a given mix of sites and diagnoses helps (Section~\ref{sec:prop2}), and which comparisons are free of site effects (Section~\ref{sec:prop3}). Section~\ref{sec:prop4} relates the answers to the scores of trained models. Proofs, two further propositions, boundary cases and a Monte Carlo check are in Supplementary Section~\ref{app:proofs}.

\subsection{Class-pure sites, designs and targets}\label{sec:setting}

We observe participants $i$ at sites $s\in\{1,\dots,S\}$. Each participant has a feature vector $x_i\in\mathbb{R}^d$, here the edges of a functional connectivity matrix, and a diagnosis $y_i\in\{0,1\}$, with $1$ for a case. The training \emph{design} is the table of how many cases $a_s$ and controls $t_s$ each site contributes, with totals $N_A=\sum_s a_s$ and $N_T=\sum_s t_s$. A site is \emph{class-pure} if it contributes only one class, that is, if $a_st_s=0$. We study training sets in which each class comes from $k$ class-pure sites with $m$ participants each, so that each class has $n=km$ participants. Such a set can be augmented with \emph{bridging subjects}: $b$ participants of the other class, recruited at the same site and processed through the same acquisition pipeline. \emph{One-way} bridging adds controls only to the case sites; \emph{two-way} bridging adds the missing class at every site (Fig.~\ref{fig:overview}).

Throughout, we use the additive model
\begin{equation}
x_i=\mu+\Delta\,y_i+u_{s(i)}+\varepsilon_i ,\label{eq:model}
\end{equation}
where $\mu$ is the overall mean, $\Delta$ is the average difference between cases and controls, $u_s$ is the shift added by site $s$, and $\varepsilon_i$ is individual variation. The $\varepsilon_i$ are i.i.d.\ with mean $0$ and covariance $\Sigma_\varepsilon$, independent of the site effects and of the design. The results below need different subsets of two assumptions.
\begin{itemize}
\item[(A0)] \emph{Additivity}: the site effect is the same for both classes within a site (no site-by-diagnosis interaction).
\item[(A1)] \emph{Exchangeable site effects}: the $u_s$ are i.i.d.\ with mean $m_u$ and covariance $\Sigma_u$, independently of which class a site contributes. Mean statements need only the weaker (A1$'$): $\mathbb{E}[u_s]$ does not depend on the class a site contributes.
\end{itemize}
In words, (A0) says that a scanner changes the features of cases and controls in the same way, and (A1) says that case sites are not systematically different from control sites. (A1$'$) fails when case sites and control sites recruit systematically different populations, and such sampling differences can be as large as disorder effects~\cite{yamashita2019}. When no site holds both classes, neither assumption can be checked from the data.

What a clinic needs is a score $f$ that ranks a case above a control \emph{at a new site}, measured by the within-site AUC, $\mathrm{AUC}_W$. To detect site shortcuts we also use the \emph{reverse set}: held-back participants at the training sites whose diagnosis is opposite to the class that their site contributed in training. A model that ranks by site puts them in the wrong order, so a reverse-set AUC far below $0.5$ signals a site shortcut. Parker and Leek used the same contrast with two genomic batches~\cite{parker2012}.

\subsection{Class-pure training is not identifiable}\label{sec:prop1}

If hospital A contributes only cases with a feature mean of $2$ and hospital B only controls with a mean of $0$, the difference may be a disease effect, a hospital effect or any mix of the two; more participants measure the two means more precisely but do not tell these explanations apart. Proposition~\ref{prop:pure} states this in general, with the error that remains under (A1).

\begin{proposition}[Class-pure sites]\label{prop:pure}
Suppose every training site is class-pure.
(a) If the site effects are unrestricted, $(\Delta,\{u_s\})$ and $(\Delta+\delta,\{u_s-\delta\,\gamma_s\})$, where $\gamma_s$ is the class of site $s$, generate the same distribution of the training data for every $\delta\in\mathbb{R}^d$; hence any estimator $T$ satisfies $\max\{\mathbb{E}\|T-\Delta\|^2,\mathbb{E}\|T-\Delta-\delta\|^2\}\ge\|\delta\|^2/4$.
(b) Under (A1), with $k$ sites of $m$ subjects per class, the difference of the two class means $\hat\Delta$ is unbiased and
\begin{equation}
\operatorname{Cov}(\hat\Delta)=\frac{2}{k}\Sigma_u+\frac{2}{km}\Sigma_\varepsilon .\label{eq:pure-cov}
\end{equation}
\end{proposition}
Part~(a) holds because $\Delta\gamma_s$ and $u_s$ enter the data only through their sum; it is the multi-site form of the observation that completely confounded batch designs admit no correction~\cite{song2020}. The site term in Eq.~\eqref{eq:pure-cov} does not shrink as $m$ grows: adding subjects to class-pure sites cannot substitute for adding sites.

\subsection{The design coefficient}\label{sec:prop2}

For any design define
\begin{equation}
c_s=\frac{a_s}{N_A}-\frac{t_s}{N_T},\qquad L=\sum_s c_s^2 .\label{eq:L}
\end{equation}

Here $c_s$ is the share of all training cases that site $s$ contributes minus its share of all training controls. $L$ is not a loss or a probability. The next result shows that $L$ is exactly the weight of the site effects in the difference between the class means.

\begin{proposition}[Design coefficient]\label{prop:L}
Under (A0), $\bar x_A-\bar x_T=\Delta+\sum_s c_su_s+(\bar\varepsilon_A-\bar\varepsilon_T)$. Under (A1) this contrast is unbiased and
\begin{equation}
\operatorname{Cov}(\bar x_A-\bar x_T)=L\,\Sigma_u+\Big(\frac1{N_A}+\frac1{N_T}\Big)\Sigma_\varepsilon .\label{eq:L-cov}
\end{equation}
$L=0$ if and only if all non-empty sites have the same case fraction; a 1:1 ratio is not required. For the symmetric designs,
\begin{equation}
L_{\mathrm{pure}}=\frac2k,\qquad
L_{\mathrm{one}}=\frac2k\Big(\frac{m}{m+b}\Big)^2,\qquad
L_{\mathrm{two}}=\frac2k\Big(\frac{m-b}{m+b}\Big)^2 .\label{eq:L-special}
\end{equation}
\end{proposition}

$L$ depends only on the design table, so it can be computed before any data are collected. For $k=2$ and $m=20$, five two-way bridges per site reduce $L$ from $1$ to $0.36$, which in the site-variance term equals the effect of increasing $k$ from $2$ to about $5.6$ ($k_{\mathrm{eff}}=2/L$). The next result says how to spend a bridging budget.

\begin{corollary}[Allocation]\label{cor:alloc}
Let $U$ controls be added to the case sites and $V$ cases to the control sites. Over real-valued allocations: (i) for fixed $U$ and $V$, $L$ is minimized by spreading each budget evenly over the sites of its side. (ii) if the total budget $B=U+V$ may be split between the sides and $B\le 2km$, $L$ is minimized by $U=V=B/2$, i.e.\ $B/(2k)$ bridges per site. The condition $B\le 2km$ cannot be dropped: for $k=1$, $m=2$, $B=5$, the split $(4,1)$ gives $L=0$ while the even split gives $L>0$.
\end{corollary}

\subsection{Within-site contrasts remove the site term}\label{sec:prop3}

Bridging reduces the site term but leaves part of it when $L>0$. A case and a control from the same site carry the same site shift $u_s$, which cancels in their difference.

\begin{proposition}[Within-site contrasts]\label{prop:within}
(a) Under (A0), for a case $a$ and a control $t$ at the same site, $x_a-x_t=\Delta+\varepsilon_a-\varepsilon_t$. Any deterministic procedure that uses only such within-site differences, with a fixed penalty and tuning rule, is invariant to arbitrary site shifts $u_s\mapsto u_s+\eta_s$.
(b) Let $h_s=a_st_s/(a_s+t_s)$, $H=\sum_sh_s$, and, at each site with $a_st_s>0$, $D_s=\bar x_{sA}-\bar x_{sT}$; suppose $H>0$, and let sums involving $D_s$ run over these sites. The least-squares slope of $x$ on $y$ with site fixed effects equals $\tilde\Delta=H^{-1}\sum_sh_sD_s$, which is unbiased with covariance $\Sigma_\varepsilon/H$.
\end{proposition}

Class-pure sites carry no weight, and the effective sample size is $H$, which is smaller than the number of within-site pairs. Three limits apply. Additivity does not imply a common linear logit (a common Gaussian within-class covariance suffices), so bridging subjects must pass through the same acquisition pipeline as the site's original class. Invariance holds in fixed feature coordinates; a global standardization fitted on the training set depends on the site composition, so the implemented pipeline is invariant only approximately. Within-site training does not remove the site constant from the scores, so pooled cross-site AUC, thresholds and calibration are still affected.

Another approach is to make the representation independent of the site, but on class-pure data this also removes the diagnostic information. If $Y=\gamma(S)$ and $P(Z\mid S=s)$ is the same for all sites, then $Z\perp Y$, and any score computed from $Z$ alone has AUC $0.5$ on the training distribution~\cite{zhao2019}; methods with a finite invariance penalty do not reach this limit, so their effect is an empirical question (Section~\ref{sec:res-pureonly}). Class-conditional alignment, $Z\perp S\mid Y$~\cite{zhao2020cfnet}, constrains only the differences among sites of the same class and leaves the difference between the average case site and the average control site, which is confounded with $\Delta$, unchanged.

\subsection{From the design to the scores}\label{sec:prop4}

To see how the site term reaches the predictions, consider a simplified linear scorer whose weight vector is a fixed matrix $Q$ times the estimated diagnostic direction, and the score gap $G$ of a new case from training site $s$ minus a new control from training site $r$. Proposition~\ref{prop:leak} in the Supplementary Material, which we derived after the experiments to explain their pattern, shows under (A0) and (A1), for new participants whose residuals are independent of the training sample given the shared site effects, that the pooled scorer has $\mathbb{E}\,G=\Delta^\top Q\Delta+(c_s-c_r)\operatorname{tr}(Q\Sigma_u)$. In the symmetric designs (with $b\le m$ under two-way bridging), a reverse-set pair has $c_s-c_r=-\sqrt{2L/k}$, so the expected shortcut term grows with $\sqrt{L}$; Section~\ref{sec:res-bridge} tests the corresponding relation between $L$ and the reverse-set AUC. For a same-site pair the site term vanishes: a score offset shared by all participants of a site cancels, so within-site ranking cannot identify such offsets of a fixed model, although the site term can still change the learned direction. The within-site scorer $Q\tilde\Delta$ has $\mathbb{E}\,G=\Delta^\top Q\Delta$ for every pair. Supplementary Section~\ref{sec:synthetic} gives a synthetic example. These statements concern mean score gaps of linear scorers with a fixed $Q$; equal mean gaps do not imply equal AUCs, and they give no guarantee for the AUC of regularized logistic regression or for how its penalty is selected.

% Section 4 -- Framework: evaluation, design, estimation (condensed 2026-11-23).
\section{A site-level framework}\label{sec:framework}

When sites are class-pure, independent evidence comes from sites: more subjects at a site only estimate that site's mean more precisely (Proposition~\ref{prop:pure}). The framework turns this into three rules, all of which work with standard models.

\subsection{Evaluation}\label{sec:eval}

\begin{table}[tbp]
\centering
\caption{The four AUCs used in this paper. All use the same scoring direction: a higher score means more case-like, and each AUC is the fraction of case and control pairs in which the case receives the higher score. An AUC of $0.12$ is therefore not an accuracy of $12\%$: the control scored higher in $88\%$ of pairs.}
\label{tab:aucs}
\makebox[\linewidth][c]{\begin{minipage}{15cm}\centering
\begin{tabular}{>{\raggedright}p{0.2\linewidth}>{\raggedright}p{0.36\linewidth}>{\raggedright\arraybackslash}p{0.32\linewidth}}
\toprule
AUC & Which cases are compared with which controls & Question it answers\\
\midrule
Subject-level cross-validation & Participants held out from the training sites; the case and control of a pair usually come from different training sites & How good does the model look on the sites it was trained on?\\
LSPO & Participants of one held-out case site against those of one held-out control site & Does the ranking carry over to sites that were not used for training?\\
Within-site AUC on new sites (main metric) & Cases and controls from the same new mixed site, averaged over test sites & Can the model separate the classes when site gives no advantage?\\
Reverse-set AUC (stress test) & Cases held out at control sites against controls held out at case sites & Does the model follow the site when the training link between site and diagnosis is reversed?\\
\bottomrule
\end{tabular}
\end{minipage}}
\end{table}

Table~\ref{tab:aucs} summarizes the four AUCs. The within-site AUC measures diagnostic ranking but cannot reveal by itself whether scores are shifted between sites; the reverse set is added for that purpose.

\paragraph{Do not use subject-level cross-validation} When folds are drawn over subjects, every test subject shares its site with training subjects of the same class, so a site shortcut is rewarded like a diagnostic signal. By Proposition~\ref{prop:pure} it persists as $m$ grows.

\paragraph{Leave-site-pair-out (LSPO)} Each fold holds out one case site and one control site, trains on the remaining $k-1$ sites of each class and computes the AUC on the two held-out sites; the estimate is the mean over all $k^2$ pairs, so it requires $k\ge2$. Subject-level cross-validation is averaged over folds in the same way, because pooling scores from different models can bias the AUC in small samples~\cite{airola2011}. LSPO never shares a site between training and test, but it measures ranking across sites: for Gaussian site effects and residuals and a fixed $w$ with $w^\top\Delta>0$ and $w^\top\Sigma_\varepsilon w>0$, a case and a control from two independent new sites are ordered correctly less often than a pair from the same new site, $\mathrm{AUC}_B\le\mathrm{AUC}_W$, because the site variance enters the comparison (Proposition~\ref{prop:targets} in the Supplementary Material). LSPO also trains on one site pair fewer than the final model, so we report its empirical bias and spread.

\paragraph{Reverse-set AUC as a stress test} The reverse set (Section~\ref{sec:setting}) asks whether a model ranks by site. In a truly class-pure collection it has to be recruited, for example by holding back some bridging subjects. An AUC well below $0.5$ means that the model ranks by site or has learned the diagnostic direction backwards; an AUC at or above $0.5$ does not prove that it ignores sites, since diagnostic signal can mask site offsets.

\subsection{Design}\label{sec:design}

\paragraph{Compute $L$ before collecting data} The design coefficient (Eq.~\ref{eq:L}) depends only on the table of participants per site and class, so recruitment plans can be compared on paper through $L/L_{\mathrm{pure}}$ and the equivalent number of sites $2/L$.

\paragraph{Bridge across classes, in both directions, evenly} By Corollary~\ref{cor:alloc}, a fixed budget is best spread evenly over the sites and, for budgets $B\le2km$, split evenly between case and control sites. Bridges must be recruited like the site's original participants: controls from a different population scanned on the same machine differ by more than the site shift, which violates (A0).

\subsection{Estimation}\label{sec:estimation}

\paragraph{Within-site pairwise logistic loss (L3)} Proposition~\ref{prop:within} suggests training only on differences between a case and a control from the same site. L3 compares such pairs and penalizes the model when the case receives the lower score. Pairing is used only during training; at test time each participant is scored individually. Let $\mathcal{P}=\{(a,t): s(a)=s(t),\ y_a=1,\ y_t=0\}$ be the within-site case-control pairs, $N_{\mathrm{pair}}=|\mathcal{P}|$, and $z$ the standardized features. L3 fits $f(x)=w^\top z$ without intercept by minimizing
\begin{equation}
\frac1{N_{\mathrm{pair}}}\sum_{(a,t)\in\mathcal{P}}\log\!\big(1+e^{-w^\top(z_a-z_t)}\big)+\frac{\|w\|^2}{2CN},\label{eq:L3}
\end{equation}
where $N$ is the number of training subjects, so that $C$ has the same scale as in standard logistic regression. Class-pure sites contribute no pairs, and the site shift cancels in every difference $z_a-z_t$. L3 coincides with the exact conditional logistic likelihood~\cite{avalos2015} only when every stratum holds one case and one control.

\paragraph{Tuning} $C$ is selected by the within-site pair AUC in an inner cross-validation on the training subjects, with $\min(5,b)$ folds stratified by site and diagnosis. Because this criterion cancels the site-specific score offsets of a fixed model, it does not guarantee tuning free of the shortcut; what guards against the shortcut is the structure of Eq.~\eqref{eq:L3} itself, which the selection rule cannot provide.

% Section 5 -- Experimental setup (condensed 2026-11-23; details in the Supplementary Material).
\section{Experimental setup}\label{sec:setup}

To measure the effect of class-pure sites against a known reference, we construct class-pure training sets from cohorts in which most sites contain both classes and evaluate on held-out mixed sites; two naturally class-pure collections provide real examples. All criteria were fixed in writing before the corresponding results were computed (Supplementary Section~\ref{app:criteria}).

\subsection{Datasets}\label{sec:data}

Table~\ref{tab:data} summarizes the three cohorts. Collections from the same institution were merged into one site, because they may share a scanner. For ABIDE-I we used the Preprocessed Connectomes Project release~\cite{craddock2013pcp} of ABIDE-I~\cite{dimartino2014} (CPAC pipeline, band-pass filtering, no global signal regression), with all $1035$ subjects that have a functional file. Each participant is represented by the Fisher-$z$ transformed Pearson correlations between the $200$ regions of the CC200 atlas~\cite{craddock2012}. ABIDE-II~\cite{dimartino2017} served as an independent replication with different subjects, preprocessing (DCAN, censoring at $0.2$\,mm framewise displacement, at least $240$\,s retained) and atlas (Gordon parcellation~\cite{gordon2016} plus subcortical regions). ADHD-200 served as a second disorder (Athena release~\cite{bellec2017}, CC200 with $190$ non-empty regions, first resting-state scan passing quality control).

\begin{table}[tbp]
\centering
\caption{Cohorts and the class-pure constructions used in this study. $n$: subjects per class in a constructed training set; $k$: class-pure sites per class; $b$: bridging subjects per site.}
\label{tab:data}
\makebox[\linewidth][c]{\begin{minipage}{15cm}\centering
\begin{tabular}{>{\raggedright}p{0.25\linewidth}>{\raggedright}p{0.2\linewidth}>{\raggedright}p{0.2\linewidth}>{\raggedright\arraybackslash}p{0.2\linewidth}}
\toprule
& ABIDE-I & ABIDE-II & ADHD-200\\
\midrule
Subjects (case/control) & 505/530 & 248/320 & 325/540\\
Features (edges) & 19,900 & 61,776 & 17,955\\
Sites used for construction & 17 & 13 & 5\\
Held-out test sites per repeat & 5 & 4 & 3\\
Constructed designs & $n=40$, $k\in\{1,2,4,6\}$ & $n=20$, $k\in\{1,2,4\}$ & $n=20$, $k=1$\\
Bridging & $b\in\{2,5,10\}$, $k\in\{2,4\}$ & $b\in\{2,5\}$, $k\in\{2,4\}$ & $b\in\{2,5\}$\\
Repeats & 50 & 50 & 60\\
Natural class-pure collection & n.a. & KUL\_3 (ASD only) & WashU (TD only)\\
\bottomrule
\end{tabular}
\end{minipage}}
\end{table}

\subsection{Constructing class-pure training sets}\label{sec:construction}

In each repeat, some sites are held out for testing. From the remaining pool, $k$ sites contribute only cases and $k$ other sites only controls, with $n$ subjects per class (design P, the pure baseline B0). The mixed design M uses the same $2k$ sites and the same number of subjects but draws both classes from every site, so that M$-$P is the cost of class purity. Participants of the opposite class who remain at the class-pure training sites form the reverse set. Bridged designs add $b$ subjects of the missing class to each case site (one-way) or to every site (two-way). The same-class control S0 adds the same number of subjects of the class that each site already contributes. Within a repeat, every method is trained on the same subjects and evaluated on the same test and reverse sets, and contrasts are paired differences. In ABIDE-I we also vary $n\in\{20,40,80\}$ at $k=2$. Design parameters were set from the site tables before any classification result on the cohort was computed. A site can take a role only if it has enough participants of that class, so the construction controls part, but not all, of the difference between case sites and control sites (Section~\ref{sec:limitations}).

For the natural collections KUL\_3 (ABIDE-II, $21$ adult male cases) and WashU (ADHD-200, $47$ controls), we compare subject-level five-fold cross-validation with leave-one-site-out evaluation on the mixed sites plus the collection. The difference between the drop in AUC for the collection and for mixed-site participants of the same class, each scored against the opposite class from the mixed sites (DiD), describes how much more the collection loses. Both AUCs pool out-of-fold scores, and the collection also differs from the mixed sites in population, so the DiD is descriptive, not causal.

\subsection{Methods compared}\label{sec:methods}

All linear models are $\ell_2$-penalized logistic regressions on standardized features. For the pooled models (B0, M, S0, C0), $C\in\{10^{-3},10^{-2},10^{-1},1\}$ is selected by stratified five-fold cross-validation over training subjects, as in common practice. L3 (Eq.~\ref{eq:L3}) uses $C\in\{10^{-3},\dots,10\}$ selected as in Section~\ref{sec:estimation}; L1 is logistic regression with one intercept per site, which absorbs additive shifts but must drop class-pure sites, whose intercepts diverge. Methods that use only class-pure data comprised tuning by LSPO, removal of the site-mean subspace, per-site standardization, Group DRO with one group per site~\cite{sagawa2020}, and three one-class models; ComBat with diagnosis as a protected covariate was applied to bridged designs, fitted once on the outer training set. We reproduced two published multi-site methods with public code: the population graph convolutional network~\cite{parisot2018}, which connects participants by sex, site and feature similarity, and the pipeline of Kunda et al.~\cite{kunda2023}, which reduces the dependence of tangent-space connectivity on site by maximum independence domain adaptation (MIDA). Each reproduced its published ABIDE-I accuracy within three percentage points. Both are transductive, so we ran them in their original setting and, through an extension of the authors' code added after the first results, with the reverse set kept out of training and scored afterwards (Supplementary Section~\ref{app:more}). Supplementary Table~\ref{tab:designs} shows how the main designs and methods differ. M versus P gives the cost of class purity, C0 versus S0 tests whether the added participants must belong to the other class, and L3 versus C0 compares two complete training procedures on the same participants, which differ in the objective and in how $C$ is selected.


\subsection{Metrics}\label{sec:metrics}

The primary metric is the within-site AUC, computed in each held-out test site and averaged across test sites; we also report the reverse-set AUC (Table~\ref{tab:aucs}). In every AUC, tied scores count one half. Each metric is computed within a repeat and then averaged across repeats. Ranges are the $2.5$th and $97.5$th percentiles across repeats; because repeats share subjects, they describe variability across overlapping splits and are not confidence intervals. For evaluation estimators, the reference is the within-site AUC of the same training design on the held-out test sites. Because repeats share participants, we applied the corrected resampled $t$-test~\cite{nadeau2003} with Holm correction within pre-specified families; its error rate under site resampling is not known, so we treat it as an approximate check (all $82$ tests in Supplementary Section~\ref{app:robust}) and rest the conclusions on replication across three cohorts with no shared participants.

% Section 6 -- Results (condensed 2026-11-23; moved material is in the Supplementary Material, Section S3).
\section{Results}\label{sec:results}

Unless stated otherwise, results are for ABIDE-I with the CC200 atlas, as means over $50$ repeats with percentile ranges in brackets and, for paired differences, the number of repeats with the expected sign.

\subsection{Class purity costs accuracy and teaches a site shortcut}\label{sec:res-problem}

With one site per class, the class-pure design P reached a within-site AUC of $0.578$ on new sites, against $0.636$ for the mixed design M trained on the same sites and the same number of participants (Fig.~\ref{fig:problem}a). The reverse-set AUC of P was $0.03$: the model ranked almost every reverse-set participant by their site (Fig.~\ref{fig:problem}b). More sites per class reduced both effects; from $k=1$ to $k=6$ the cost M$-$P fell from $0.057$ to $0.015$ and the reverse-set AUC rose to $0.36$. Adding participants to the same sites did not help and made the shortcut stronger: at $k=2$, quadrupling $n$ from $20$ to $80$ raised the within-site AUC of P by only $0.034$, against $0.080$ for M, and the reverse-set AUC of P fell from $0.16$ to $0.08$ (Fig.~\ref{fig:problem}c), as predicted by the site term $(2/k)\Sigma_u$ of Proposition~\ref{prop:pure}.

The shortcut is available because sites are easy to recognise: a multinomial logistic regression identified the site with a balanced accuracy of $0.72$ to $0.92$ in the three cohorts (chance at most $0.17$), also within one diagnosis (Supplementary Fig.~\ref{fig:sitecode}).

\begin{figure}[tbp]
\centering
\widefig{figures/fig_problem.pdf}
\caption{Class purity costs accuracy on new sites and teaches a site shortcut that more sites reduce and more subjects do not (ABIDE-I, CC200). (a)~Within-site AUC on held-out test sites for the class-pure design P (red) and the mixed design M (grey), trained on the same sites with $n=40$ subjects per class; numbers at the top are the mean cost M$-$P. (b)~Reverse-set AUC of P; values in the shaded region below chance mean that the model ranks reverse-set subjects by their site. (c)~At $k=2$, more subjects per class raise the within-site AUC of M more than that of P and lower the reverse-set AUC of P (hollow markers, dashed line). Small dots are single repeats ($50$ per condition), large markers are means.}
\label{fig:problem}
\end{figure}

\subsection{Subject-level cross-validation overstates class-pure models}\label{sec:res-eval}

At $k=2$, subject-level cross-validation estimated a within-site AUC of $0.94$ for a model that reached $0.575$ on new sites, an overestimate of $0.37$ in all $50$ repeats; it shrank with $k$ but remained $0.19$ at $k=6$ (Fig.~\ref{fig:evaluation}a--c). LSPO had a mean error of $-0.05$ at $k=2$ and below $0.01$ in magnitude at $k=4$ and $6$, and its mean absolute error was less than half that of subject-level cross-validation at every $k$. Its spread remained large (mean absolute error $0.09$--$0.16$), and neither estimator tracked the test-site value from one repeat to the next; for a single study, LSPO is a rough estimate.

The natural collections showed the same optimism (Fig.~\ref{fig:evaluation}d). The descriptive DiD (Section~\ref{sec:construction}) was, for the case-only KUL\_3, $+0.30$ [$0.24$, $0.35$], an exploratory AUC version of a pre-specified analysis of sensitivity; for the control-only WashU it was $+0.30$ [$0.28$, $0.34$], pre-specified; and for Pittsburgh ($92$ controls, $4$ cases) $+0.50$ (Supplementary Section~\ref{app:more}).

\begin{figure}[tbp]
\centering
\widefig{figures/fig_evaluation.pdf}
\caption{Subject-level cross-validation overstates class-pure models, and leave-site-pair-out does not. (a--c)~Estimates of the within-site AUC of P from the training data against the value reached on held-out test sites (ABIDE-I, CC200, $n=40$): subject-level five-fold cross-validation (red) and LSPO (blue). Dots are single repeats ($50$ per panel), large markers are means, the dashed line is $y=x$, and the mean signed error is given above each panel. (d)~Naturally class-pure collections: AUC under subject-level cross-validation (open marker) and leave-one-site-out evaluation (filled marker) for the collection (red) and for mixed-site subjects of the same class (grey); DiD is the difference between the two drops. KUL\_3 (ABIDE-II) contains only cases; WashU only controls; Pittsburgh $92$ controls and $4$ cases (ADHD-200).}
\label{fig:evaluation}
\end{figure}

\subsection{Methods without bridges, including two published ones, keep the shortcut}\label{sec:res-pureonly}

The methods we tested on class-pure data did not materially reduce the measured shortcut. Tuning by LSPO, projecting out the site-mean subspace and Group DRO with one group per site changed the reverse-set AUC by less than $0.005$. Per-site standardization gave an AUC of exactly $0.5$, since subtracting the site mean also subtracts the class mean. One-class models avoided the shortcut only by learning no diagnostic signal (within-site AUC $0.50$--$0.55$) and met neither of the predictions we had made for them (Supplementary Section~\ref{app:more}).

Two published multi-site methods failed in the two ways that Section~\ref{sec:theory} describes (Supplementary Fig.~\ref{fig:published}). With the reverse set kept out of training, as in a real class-pure collection, the population graph network~\cite{parisot2018} showed the strongest shortcut of all models (reverse-set AUC $0.02$ at $k=2$ and $0.22$ at $k=4$, below $0.5$ in all $50$ repeats), because its graph links participants of the same site; subject-level cross-validation rated it at $1.00$ against $0.56$ on new sites. The MIDA pipeline~\cite{kunda2023} instead lost most of its diagnostic signal (within-site AUC $0.52$ and $0.53$, against $0.61$ with mixed training sites). Two-way bridging raised the reverse-set AUC over B0 by $0.29$ to $0.34$ for the graph network and $0.08$ to $0.18$ for MIDA.


\subsection{Bridging works through cross-class contrast, and $L$ ranks designs before training}\label{sec:res-bridge}

A few participants of the missing class at each site recovered much of the lost accuracy. With two-way bridging and $b=5$ at $k=2$, the pooled model C0 improved the within-site AUC from $0.575$ to $0.614$ and the reverse-set AUC from $0.12$ to $0.30$; with $b=10$ at $k=4$ the reverse-set AUC reached $0.68$. The added participants had to belong to the other class: against the same-class control S0, which adds the same number of participants of each site's own class, C0 was higher by $0.015$--$0.044$ in within-site AUC and $0.10$--$0.29$ in reverse-set AUC across the four $b=5$ cells, based on the $15$ to $35$ repeats in which both designs were feasible (Fig.~\ref{fig:bridging}a,b).

The design coefficient, which needs only the planned site-by-class table (Proposition~\ref{prop:L}), predicted these differences before any model was trained: the larger the $L$ of a design, the lower its reverse-set AUC (Fig.~\ref{fig:bridging}d). The rank correlation between $L$ and the reverse-set AUC of C0 was $-0.997$ with CC200 and $-0.991$ with the Harvard-Oxford atlas, but most of these cells had been inspected before Proposition~\ref{prop:L} was written, so this correlation is exploratory. Two tests were specified in advance. Spreading a fixed budget evenly over sites, as Corollary~\ref{cor:alloc} prescribes, gave a higher reverse-set AUC than concentrating it on half of the sites, by $0.040$ (CC200) and $0.037$ (Harvard-Oxford) at $k=4$ (Fig.~\ref{fig:bridging}c), and on ABIDE-II the rank correlation over $11$ new cells was $-0.970$.

ComBat fitted on the bridged training set reached within-site AUCs within $0.01$ of C0, but its reverse-set AUC failed a pre-specified null test.

\begin{figure}[tbp]
\centering
\widefig{figures/fig_bridging.pdf}
\caption{Cross-class bridges outperform same-class additions, even allocation works best, and $L$ ranks designs before training (ABIDE-I; panel d also ABIDE-II). (a, b)~Is crossing needed? Paired differences between the pooled model with cross-class bridges (C0) and with the same number of added subjects of each site's own class (S0), in within-site AUC and in reverse-set AUC; brackets give the repeats in which both designs were feasible. (c)~How to spend a budget? Reverse-set AUC of C0 at $b=5$ with the budget spread evenly minus concentrated on half of the sites, on the reverse-set subjects common to both designs; circles CC200, squares Harvard-Oxford, light blue one-way and dark blue two-way bridging. (d)~Does $L$ rank designs? $L$ against the mean reverse-set AUC of C0 in every design cell; a larger $L$ goes with a lower reverse-set AUC, that is, a stronger shortcut; shade encodes the design, marker shape the cohort and atlas, hollow markers the concentrated cells. In (a--c), dots are single repeats and large markers are means; in (d), markers are cell means.}
\label{fig:bridging}
\end{figure}

\subsection{Within-site estimation keeps the reverse set above chance}\label{sec:res-within}

L3 trains one linear scorer for all sites, but its loss only compares a case with a control from the same site; class-pure sites contribute no such pairs, and at test time every participant at a new site is scored individually (Section~\ref{sec:estimation}). Training only on these within-site contrasts removed most of what remained of the shortcut. Across the $12$ bridged cells, the reverse-set AUC of L3 lay between $0.59$ and $0.67$, whereas that of C0 ranged from $0.15$ to $0.68$ and depended strongly on $L$ (Supplementary Fig.~\ref{fig:tradeoff}). L3 had the higher reverse-set AUC in all $50$ repeats of $11$ of the $12$ cells; in the exception, $k=4$, $b=10$ two-way, $L$ is close to zero and the two estimators performed alike. With two-way bridging, the within-site AUC of L3 differed from that of C0 by $-0.003$ to $+0.009$. With one-way bridging and few bridges, L3 paid a price ($-0.029$ at $k=2$, $b=2$), because only half of the sites contribute pairs.

The protection came from the objective itself: mixing a weight of only $0.1$ of the pooled loss into L3 lowered the reverse-set AUC from $0.60$ to $0.33$, and selecting the weight by within-site AUC chose the pure within-site objective in only $41\%$ of repeats (Supplementary Fig.~\ref{fig:lambda}). This criterion cannot see the site offsets of a fixed model, which cancel within each site (Proposition~\ref{prop:leak}(c)), so it does not guard against the shortcut by itself. A linear L3 head on a frozen self-supervised encoder reached a reverse-set AUC of $0.59$ against at most $0.39$ for deep baselines; one of four pre-specified cell-wise criteria was not met (Supplementary Section~\ref{app:deep}).


\subsection{The findings replicate on a second cohort and a second disorder}\label{sec:res-replication}

On ABIDE-II, the four assessable pre-specified criteria were met (Fig.~\ref{fig:replication}; Supplementary Table~\ref{tab:replication}): class purity cost $0.050$ at the smallest $k$ and nothing with more sites, subject-level cross-validation overstated by $0.38$, the rank correlation of $L$ with the reverse-set AUC of C0 was $-0.970$ over $11$ new cells, and L3 reached a reverse-set AUC of $0.63$ against $0.39$ for C0 at a test cost of $0.006$; the same-class control could not be assessed because fewer than $20$ repeats were feasible. On ADHD-200, where only $k=1$ could be constructed, the pooled model ranked the reverse set almost entirely by site (AUC $0.08$) and L3 restored it ($0.52$) at equal test AUC, but the cost of purity did not meet its threshold; even the full training pool reached a within-site AUC of only $0.513$, so this cohort carries little diagnostic signal in connectivity, and L3 removed the shortcut without learning more about ADHD.

\begin{figure}[tbp]
\centering
\widefig{figures/fig_replication.pdf}
\caption{The three main findings replicate on ABIDE-II (Gordon atlas, $n=20$; top row) and on ADHD-200 (CC200, $n=20$, only $k=1$; bottom row). (a, d)~Within-site AUC of P (red) and M (grey) and reverse-set AUC of P (hollow markers, dashed line). (b, e)~Estimates from training data against the within-site AUC on new sites, for subject-level cross-validation (red) and LSPO (blue), which requires $k\ge2$. (c, f)~Reverse-set AUC of C0 (red) and L3 (blue) in each bridged cell; no arrow is drawn where the two differ by less than $0.04$. Dots are single repeats, large markers are means, and dashed lines mark chance or $y=x$.}
\label{fig:replication}
\end{figure}

The pre-specified criteria for the cost of purity, the bias of subject-level cross-validation and the cross-class effect of bridging were also met on the $873$ ABIDE-I participants passing all quality ratings, with the Harvard-Oxford atlas and with BrainNetCNN~\cite{kawahara2017}, although the BrainNetCNN cost only just exceeded its threshold (Supplementary Section~\ref{app:robust}).


\subsection{The rules carry over to partially confounded sites}\label{sec:res-partial}

Most multi-site data lie between mixed and class-pure sites: every site holds both diagnoses, in different proportions. We varied this continuously at fixed size and prevalence, from balanced sites to design P (Fig.~\ref{fig:partial}; Supplementary Section~\ref{app:partial}). The reverse-set AUC of the pooled model fell steadily with $L$, from $0.67$ to $0.48$ at $L/L_{\mathrm{pure}}=0.16$ and $0.28$ at $0.64$, below its balanced value in all $50$ repeats at every level, and subject-level cross-validation overstated the model in step (by $0.06$ and $0.22$), whereas leaving one training site out erred by $0.01$ to $0.04$. Site-by-class weights and ComBat kept most of the shortcut. L3 kept the reverse-set AUC between $0.60$ and $0.66$, above the pooled model in all $50$ repeats from $L/L_{\mathrm{pure}}=0.16$, at a within-site AUC $0.001$ to $0.007$ higher up to $0.36$ and $0.017$ lower at $0.64$, where each minority cell holds two participants. One intercept per site (L1) performed like L3, as on the bridged designs (Supplementary Section~\ref{app:criteria}).

\begin{figure}[tbp]
\centering
\widefig{figures/fig_partial.pdf}
\caption{Under partial confounding, the shortcut and the bias of subject-level cross-validation grow with $L$, and within-site training removes the shortcut (ABIDE-I, CC200). Six training sites with $20$ participants each; at $3$ case-leaning sites a fraction $0.5+\delta$ are cases, at $3$ control-leaning sites the same fraction are controls, so that $L/L_{\mathrm{pure}}=(2\delta)^2$ and $L/L_{\mathrm{pure}}=1$ is design P. The reverse set holds left-out participants of each site's minority class. (a)~Reverse-set AUC and (b)~within-site AUC on the $11$ held-out sites for the pooled model, the pooled model with site-by-class weights, ComBat with diagnosis protected before the pooled model, L3, and logistic regression with one intercept per site (L1, dashed; it coincides with L3); L1 and L3 are undefined at $L/L_{\mathrm{pure}}=1$. (c)~Error of two estimates of the within-site AUC of the pooled model from its training data; leaving one training site out is undefined at $L/L_{\mathrm{pure}}=1$. Lines join means over $50$ repeats; dots in (c) are single repeats.}
\label{fig:partial}
\end{figure}

% Section 7 -- Practical recommendations (added 2026-11-24 at the user's request: what readers should do with their own data).
\section{Practical recommendations}\label{sec:recommendations}

Table~\ref{tab:recommend} turns the results into recommendations for three situations.

\begin{table}[tbp]
\centering
\caption{What to do with class-pure sites. In every study, also report the table of subjects by site and diagnosis, site-level performance and the reverse-set AUC (Section~\ref{sec:eval}). Section, proposition and corollary numbers give the supporting result.}
\label{tab:recommend}
\makebox[\linewidth][c]{\begin{minipage}{15cm}\centering\small
\begin{tabular}{>{\raggedright}p{0.2\linewidth}>{\raggedright\arraybackslash}p{0.74\linewidth}}
\toprule
Situation & What to do\\
\midrule
Planning a new multi-site study &
$\bullet$ Compute $L$ (Eq.~\ref{eq:L}) from the planned subjects per site and diagnosis for each candidate plan, and choose the plan with the smallest $L$ (Sections~\ref{sec:prop2} and~\ref{sec:res-bridge}).\newline
$\bullet$ Recruit a few participants of the missing diagnosis at every class-pure site, spread evenly over sites (Corollary~\ref{cor:alloc}) and recruited and scanned like the site's other participants (Section~\ref{sec:design}).\\
\addlinespace
Existing data with class-pure sites, no bridges &
$\bullet$ Evaluate by leave-site-pair-out, not subject-level cross-validation; this needs two or more sites per diagnosis (Sections~\ref{sec:eval} and~\ref{sec:res-eval}).\newline
$\bullet$ Add sites rather than participants at the same sites (Section~\ref{sec:res-problem}).\newline
$\bullet$ Harmonization, reweighting or site invariance alone do not remove the shortcut (Section~\ref{sec:res-pureonly}); state that the claim assumes exchangeable sites (Proposition~\ref{prop:pure}).\\
\addlinespace
Some sites hold both diagnoses &
$\bullet$ Train with L3; class-pure sites drop out of the loss (Section~\ref{sec:res-within}). The larger $L$ of the training data, the more a pooled model can rely on the site (Section~\ref{sec:res-bridge}).\\
\bottomrule
\end{tabular}
\end{minipage}}
\end{table}

% Section 8 -- Limitations (shortened to the two main limitations, 2026-11-24, user decision).
\section{Limitations}\label{sec:limitations}

The main limitation is that the class-pure sites were constructed from cohorts in which the sites hold both classes. Drawing both roles at random from the same cohort controls part, but not all, of the difference between case sites and control sites, and does not guarantee (A1$'$), because a site can take a role only if it has enough participants of that class. In real collections, cases from a clinic and controls from a school may also differ in age, medication or motion, a bias that none of our methods removes; the natural collections show that the evaluation bias is real, but not whether bridging removes it there. The constructed designs were also small, so LSPO remained imprecise (mean absolute error $0.09$ to $0.19$ across the two ABIDE cohorts).

The theory assumes additive site effects (A0) and linear scores. Bridging and within-site training remove a site shift common to both classes but cannot address a site-by-diagnosis interaction, such as a scanner that distorts patients' signals differently. Proposition~\ref{prop:leak} was formulated after the experiments, and our L3 implementation is shift-invariant only approximately.

% Section 9 -- Conclusions. PR: longer than the abstract; strengths, weaknesses, benefit to others, future work.
\section{Conclusions}\label{sec:conclusions}

When each site contributes a single diagnosis, subject-level cross-validation rewards models that score by site, and without an assumption that the data cannot test, diagnosis and site cannot be separated at all. Holding out whole sites removed an optimism of $0.37$ to $0.46$ in AUC; with bridging subjects, what mattered was that they crossed the class boundary, and $L$ ranked designs before training; training only on within-site comparisons kept the reverse-set AUC above chance, with test AUC within $0.01$ of the pooled model when every site was bridged. When every site held both diagnoses in unequal proportions, the shortcut and the optimism of subject-level cross-validation grew with $L$, and within-site training again removed the shortcut.

The main strength of the framework is that it acts on the study design and the training objective, so it can be planned before data collection and combined with standard classifiers, and each rule was tested with pre-specified criteria on three cohorts covering two disorders that share no participants. Its main weaknesses are that class purity was constructed rather than observed and that the theory covers additive site effects and linear scores (Section~\ref{sec:limitations}).

For researchers who pool clinical data, Table~\ref{tab:recommend} lists what to do in each situation; the two simplest steps, tabulating subjects by site and diagnosis and reporting site-level performance, apply to every study, and a few bridging participants per site are worth more than as many additions of each site's own diagnosis; with within-site training two per site were enough, whereas pooled training needed bridges to make up a third of each site. For spurious correlations without bias-conflicting samples in pattern recognition, our results point to a remedy at the level of the data: collecting a few such samples by design. Future work should recruit bridging subjects prospectively at class-pure clinical sites, extend the theory to site-by-diagnosis interactions and non-linear scores, and test within-site objectives for EEG, structural MRI and end-to-end deep models.


\section*{CRediT authorship contribution statement}
\textbf{Jincheng Hu:} Conceptualization, Methodology, Software, Formal analysis, Investigation, Visualization, Writing -- original draft. \textbf{Ping Xie:} Conceptualization, Funding acquisition, Supervision, Writing -- review \& editing. \textbf{Sha Wang:} Data curation, Visualization, Writing -- review \& editing. \textbf{Jiaxin Li:} Data curation, Visualization, Writing -- review \& editing. \textbf{Guoqian Jiang:} Conceptualization, Formal analysis, Supervision, Writing -- review \& editing.
\section*{Data availability}
ABIDE-I, ABIDE-II and ADHD-200 are public. Code and summary tables are available at \url{https://github.com/123jc-hu/class-pure-sites}.
\section*{Declaration of generative AI and AI-assisted technologies\\ in the manuscript preparation process}
During the preparation of this work the authors used Claude (Anthropic) to improve the language and readability of the manuscript. After using this tool, the authors reviewed and edited the content as needed and take full responsibility for the content of the published article.
\section*{Acknowledgements}
This work was supported by the National Natural Science Foundation of China (62471429).


% let reference entries break across pages; otherwise unbreakable entries leave gaps at page bottoms
\makeatletter\patchcmd{\thebibliography}{\clubpenalty4000}{\clubpenalty0}{}{\typeout{PATCHFAIL club}}\patchcmd{\thebibliography}{\widowpenalty4000}{\widowpenalty0}{}{\typeout{PATCHFAIL widow}}\makeatother
\bibliographystyle{elsarticle-num}
\bibliography{refs}

\end{document}
